% A technical note on how the Realistic Planets mod's map of Neptune is made. build.sh makes its
% figures and compiles it with Tectonic into the mod's folder, after ../neptune/build.sh.
\documentclass[11pt,a4paper]{article}
\usepackage[british]{babel}
\usepackage[a4paper,margin=2.4cm]{geometry}
\usepackage{amsmath,amssymb}
\usepackage{fontspec}
\usepackage{graphicx}
\usepackage{booktabs}
\usepackage{siunitx}
\usepackage[round,authoryear]{natbib}
\usepackage{xcolor}
\usepackage{pgfplots}
\pgfplotsset{compat=1.18}
\usepackage[font=small,labelfont=bf]{caption}
\usepackage{microtype}
\usepackage[colorlinks=true,linkcolor=black,citecolor=black,urlcolor=blue!50!black]{hyperref}

\renewcommand{\topfraction}{0.9}
\renewcommand{\bottomfraction}{0.8}
\renewcommand{\textfraction}{0.08}
\renewcommand{\floatpagefraction}{0.8}
\setcounter{topnumber}{3}
\setcounter{totalnumber}{4}

\definecolor{dayone}{HTML}{2B6CB0}
\definecolor{daytwo}{HTML}{C0392B}
\sisetup{group-separator={,}, group-minimum-digits=4}
\newcommand{\vect}[1]{\mathbf{#1}}
\newcommand{\unitv}[1]{\hat{\mathbf{#1}}}

\hypersetup{pdftitle={A reproducible global map of Neptune from Voyager 2 narrow-angle camera images},
  pdfauthor={Kenny Van de Maele}}

\title{A reproducible global map of Neptune from Voyager~2 narrow-angle camera images}
\author{Kenny Van de Maele\\[2pt]
  \small OpenReliant\\[6pt]
  \small Technical note on the Neptune texture of the Realistic Planets mod for OpenReliant}
\date{October 2026}

\begin{document}
\maketitle

\begin{abstract}
\noindent
We describe the construction of a global map of the albedo and colour of Neptune from 34
calibrated images taken by the narrow-angle camera of Voyager~2 on 17 and 20 August 1989. The map
was made as a texture for the Realistic Planets modification of OpenReliant, an open-source
reimplementation of the video game StarLancer (1999), but is intended to be of use beyond it. As no
reconstructed pointing of the cameras was available to us, the pointing of each image was determined
from the image itself: the position and scale of the disk from a robust fit to the sunlit limb, with
a median residual of 0.47~pixels, and the roll angle from the requirement that the cloud bands
follow parallels of latitude, pooled over the images of each day, which gives \ang{-6.85} and
\ang{-6.31} with bootstrap standard deviations of \ang{0.47} and \ang{0.25}. The images were
corrected for limb darkening with Minnaert's law and for fixed-pattern blemishes of the camera,
advected to a common epoch with the zonal wind profile of \citet{sromovsky1993}, cross-calibrated
photometrically through their overlaps, and mosaicked on a grid of \ang{0.1}. The colour was
derived from the geometric albedo spectrum of \citet{karkoschka1994}, modulated by mosaics made in
the orange, green and blue filters, and converted to sRGB with the CIE~1931 colour-matching
functions; the mean colour of the disk is sRGB (138, 182, 202). The map contains data over
\SI{85}{\percent} of the surface, from the south pole to about \ang{40}N, at a resolution of 62 to
70~km per pixel; the remaining northern latitudes are filled with mirrored southern bands. The
individual methods are well established, and Voyager mosaics of Neptune at higher resolution exist;
the contribution of this note is a documented, reproducible processing pipeline and a map released
under an open licence.

\medskip\noindent\textbf{Keywords:} Neptune; Voyager~2; image mosaicking; planetary photometry;
true colour
\end{abstract}

\section{Introduction}

StarLancer (1999) represents Neptune by a sphere of 720 triangles, or 320 at large distances, with
a texture of $256\times256$ pixels. The Realistic Planets modification for OpenReliant replaces it
with a sphere of \num{16128} triangles and a map derived from spacecraft images. The map has to
remain sharp at large angular sizes: at one navigation point of the game the viewer is 0.59~million
game units above the cloud tops, where the planet subtends about \ang{86}.

Neptune has been mapped from Voyager~2 images since the encounter. The first results of the imaging
team described the planet's bands and discrete features \citep{smith1989}, and the motions of its
clouds were measured in detail by \citet{limaye1991} and \citet{sromovsky1993}. The mosaics
produced by B.~Jónsson in 2014 from about 20 green- and clear-filter images taken at ranges of 3 to
4~million~km are among the most detailed global views of the planet made since
\citep{lakdawalla2014}; no licence for their reuse is stated. The approach described here follows
the same general lines: calibrated, geometrically corrected images are projected onto a
latitude-longitude grid, with the illumination removed by a photometric model.
\citet{irwin2024} used spectra from the Hubble Space Telescope and the Very Large Telescope to
rebalance colour images from Voyager~2 and Hubble, and showed that the true colour of Neptune is a
pale greenish blue, only slightly bluer than that of Uranus. Hubble's Outer Planet Atmospheres
Legacy (OPAL) programme maps Neptune regularly at two pixels per degree
\citep{simon2015,opal}.

For use in the game, the map was required to be global, sharp at large angular sizes, close to
true colour, and available under an open licence. None of the products known to us met all four
requirements: the Hubble maps resolve about 1000~km on Neptune, no licence is stated for Jónsson's
mosaics, and the openly licensed maps of Neptune that we found are partly painted. We therefore
made a new map from the calibrated Voyager~2 images, which are in the public domain.

Section~\ref{sec:data} describes the data, Section~\ref{sec:methods} the methods,
Section~\ref{sec:results} the results, and Section~\ref{sec:game} the adjustments made for the
game. Section~\ref{sec:discussion} discusses the limitations of the map, and
Section~\ref{sec:conclusions} concludes.

\section{Data}\label{sec:data}

\subsection{Voyager 2 images}

The images were obtained from the PDS Ring-Moon Systems Node through its search service, OPUS,
which also provides the range, sub-spacecraft point and sub-solar point of each image. OPUS lists
716 narrow-angle camera images of Neptune taken between 14 and 25 August 1989. We used the GEOMED
products of the data set VG2-N-ISS-2/3/4/6-PROCESSED-V1.0, which are radiometrically calibrated to
$I/F$ \citep[the calibration of the cameras is described by][]{danielson1981} and geometrically
corrected by resampling onto $1000\times1000$ pixels of \SI{7.842}{\micro\radian}.

Two sets of images were selected (Table~\ref{tab:frames}). The luminance set consists of 10 clear-
and green-filter images taken on 20 August, which cover all longitudes at 62 to 70~km per pixel
at the centre of the disk. The colour set consists of eight triplets of orange-, green- and
blue-filter images taken on 17 August at 97 to 109~km per pixel. In all 34 images the phase angle
is \ang{15.0}, the planetocentric latitude of the sub-spacecraft point is \ang{-28.5} and that of
the sub-solar point is \ang{-23.2}, so that the southern hemisphere is well observed, while high
northern latitudes are seen obliquely or not at all. The whole disk fits within the frame at ranges
beyond about 6.3~million~km; the images of 20 August are among the last in which it does, which
matters for the determination of the pointing (Section~\ref{sec:pointing}).

\begin{table}[htbp]
\centering
\caption{The images used. Times are UTC; the resolution is that of a raw pixel at the centre of
the disk. The individual images and their geometry are listed in \texttt{frames.json} in the
sources of the map.}
\label{tab:frames}
\begin{tabular}{llll}
\toprule
Set & Date (1989) and time & km per pixel & Images \\
\midrule
Luminance & 20 August, 00:31 to 15:13 & 62 to 70 & 7 clear, 3 green \\
Colour & 17 August, 02:35 to 23:15 & 97 to 109 & 8 orange, 8 green, 8 blue \\
\bottomrule
\end{tabular}
\end{table}

\subsection{Reference spectrum}

The colour of the map is based on the geometric albedo spectrum of Neptune measured by
\citet{karkoschka1994} at the European Southern Observatory from 23 to 26 July 1993, which covers
300 to 1000~nm sampled at intervals of 0.4~nm. It is archived in the data set
ESO-J/S/N/U-SPECTROPHOTOMETER-4-V2.0 at the PDS Atmospheres Node \citep{karkoschka1998}.

\subsection{Zonal winds}

The clouds of Neptune move rapidly relative to the planet's rotation (System~III), at about
\SI{400}{\metre\per\second} westward at the equator. We used the symmetric polynomial fitted by
\citet{sromovsky1993} to the cloud motions measured in Voyager images,
\begin{equation}
U(\varphi) = -398 + 0.188\,\varphi^{2} - 1.2\times10^{-5}\,\varphi^{4}\quad\si{\metre\per\second},
\label{eq:wind}
\end{equation}
where $\varphi$ is the planetographic latitude in degrees and $U$ is positive eastward.

\section{Methods}\label{sec:methods}

\subsection{Geometric model}

Neptune is modelled as an oblate spheroid,
\begin{equation}
\frac{x^{2}+y^{2}}{a^{2}} + \frac{z^{2}}{c^{2}} = 1,
\qquad a = \SI{24764}{\kilo\metre},\quad c = \SI{24341}{\kilo\metre},
\label{eq:spheroid}
\end{equation}
in body-fixed coordinates with $z$ along the spin axis \citep{archinal2018}. Let $\vect{O}$ be the
position of the spacecraft, $\unitv{f}$ the unit vector from the spacecraft towards the centre of
Neptune, $\unitv{n}$ the unit vector along the projection of the spin axis onto the plane normal
to $\unitv{f}$, and $\unitv{r} = \unitv{f}\times\unitv{n}$. A point $\vect{P}$ on the planet, with
$\vect{d} = \vect{P} - \vect{O}$, is imaged at sample $u$ and line $v$ (increasing downwards):
\begin{gather}
\xi = \frac{\vect{d}\cdot\unitv{r}}{\vect{d}\cdot\unitv{f}}, \qquad
\eta = \frac{\vect{d}\cdot\unitv{n}}{\vect{d}\cdot\unitv{f}},
\label{eq:tangent}\\
u = u_{0} + \frac{s}{p}\left(\xi\cos\theta - \eta\sin\theta\right), \qquad
v = v_{0} - \frac{s}{p}\left(\xi\sin\theta + \eta\cos\theta\right).
\label{eq:camera}
\end{gather}
Here $p = \SI{7.842}{\micro\radian}$ is the nominal pixel scale of the GEOMED products, and four
parameters are fitted to each image: the position $(u_{0}, v_{0})$ of the centre of the planet, a
scale factor $s$ close to unity, and the roll angle $\theta$, by which the projected spin axis is
rotated counter-clockwise from the upward direction of the image. The inverse mapping intersects
the ray through each pixel with the spheroid. The cosines of the emission and incidence angles,
$\mu$ and $\mu_{0}$, are computed from the normal of the spheroid.

\subsection{Disk centre and scale}\label{sec:pointing}

Reconstructed pointing of the Voyager~2 scan platform during the Neptune encounter was not
available to us: the C-kernel distributed by NASA's Navigation and Ancillary Information Facility
(NAIF) describes the attitude of the spacecraft bus, not of the scan platform that carries the
cameras. The pointing was therefore determined from each image. An initial estimate of
$(u_{0}, v_{0}, s)$ was obtained by fitting a model disk with Minnaert limb darkening, including the
terminator, to the image for a coarse set of roll angles. This estimate proved to be biased by up
to 26~pixels, because the model is least accurate near the limb.

The final estimate uses the sunlit limb alone. Along radial lines at intervals of \ang{1}, where
the Sun is at least \ang{8.6} above the horizon ($\mu_{0} > 0.15$) at \SI{97}{\percent} of the
model limb radius, the brightness profile is sampled every 0.25~pixels; the limb point
$\vect{q}_{i}$ is located, by linear interpolation, where the profile falls to the mean of the
median brightness just inside the limb (0.90 to 0.96 limb radii) and that of the sky (beyond 1.06
limb radii). The parameters are then found by minimising
\begin{equation}
\sum_{i} \rho\!\left(\lVert\vect{q}_{i} - \vect{c}\rVert - R(\phi_{i};\, \vect{c}, s, \theta)\right),
\qquad \rho(z) = 2\left(\sqrt{1 + z^{2}} - 1\right),
\label{eq:limb}
\end{equation}
where $\vect{c} = (u_{0}, v_{0})$, $\phi_{i}$ is the direction of $\vect{q}_{i}$ from $\vect{c}$,
$R$ is the distance of the model limb from $\vect{c}$ in that direction, in pixels, and $\rho$ is a
robust loss with a scale of one pixel. The limb points are detected again with the updated model,
three times in all.

\subsection{Roll angle}\label{sec:roll}

The outline of the disk constrains the roll angle only weakly: at a phase angle of \ang{15} it is
nearly circular, and the quality of the limb fit changes by less than \SI{1}{\percent} over a range
of \ang{30} in $\theta$. The roll angle was instead determined from the cloud bands, which on
Neptune follow parallels of latitude. For a trial angle $\theta$, the image is projected onto a
grid of \ang{0.5} in latitude and longitude between \ang{-72} and \ang{30} latitude, using only
cells with $\mu, \mu_{0} > 0.45$, and the albedo (Equation~\ref{eq:albedo}) is divided by its
median $\tilde{A}$ over the image. The cost
\begin{equation}
C(\theta) = \frac{\sum_{j} n_{j}\, M_{j}(\theta)^{2}}{\sum_{j} n_{j}}, \qquad
M_{j} = \operatorname*{median}_{k} \left\lvert \frac{A_{jk}}{\tilde{A}} -
\operatorname*{median}_{k} \frac{A_{jk}}{\tilde{A}} \right\rvert,
\label{eq:cost}
\end{equation}
is the mean squared median absolute deviation of the albedo along the latitude rows $j$, each
weighted by its number $n_{j}$ of valid cells $k$; rows with fewer than 60 valid cells are excluded.
The median absolute deviation was preferred to the variance, which is dominated by discrete
features such as the Great Dark Spot.

The minima of single images scatter by several degrees. Since the roll of the camera relative to
Neptune is expected to change little within a day, the curves were evaluated from \ang{-25} to
\ang{10} in steps of \ang{0.5}, divided by their means, and summed over the images of each day. The
roll angle of the day is the vertex of the parabola through the lowest sample of the summed curve
and its two neighbours. Its statistical uncertainty was estimated by a bootstrap over the images of
each day, with 2000 resamplings. The limb fit (Equation~\ref{eq:limb}) was then repeated with the
roll angle of the day.

\subsection{Fixed-pattern blemishes}

The first mosaics showed small dark rings that remain at the same pixel positions in every image,
whereas the clouds do not. To isolate them, the band-pass-filtered brightness of each image (a
Gaussian of 1.5~pixels minus one of 10~pixels, divided by the latter) was averaged in camera
coordinates over the disks of all 34 images (Figure~\ref{fig:blemishes}). Five compact rings were
identified with a Laplacian-of-Gaussian filter and confirmed visually. Within 26~pixels of the
centre of each, every image is divided by one plus this averaged pattern. Masking the rings instead
left gaps in the mosaic where only one image observed a region well. Isolated deviant pixels,
differing from the median of their $3\times3$ neighbourhood by more than 0.02 in $I/F$, are
replaced by that median.

\subsection{Photometric model}

Limb darkening is described by Minnaert's law \citep{minnaert1941},
\begin{equation}
\frac{I}{F} = B\,\mu_{0}^{k}\,\mu^{k-1},
\label{eq:minnaert}
\end{equation}
fitted to each image by linear least squares in the form
$\ln(\mu\,I/F) = \ln B + k\ln(\mu\mu_{0})$, using pixels with $\mu, \mu_{0} > 0.25$. The albedo of
a point is then
\begin{equation}
A = \frac{I/F}{\mu_{0}^{k}\,\mu^{k-1}}.
\label{eq:albedo}
\end{equation}
For the luminance set, $A$ is divided by the fitted $B$ of each image, so that clear- and
green-filter images can be combined; for the colour sets, $B$ is retained, since the ratios between
the filters carry the colour.

\subsection{Photometric cross-calibration}

Equation~\eqref{eq:minnaert} is least accurate at large emission and incidence angles, and north of
the equator all images observe the clouds obliquely. The first mosaics therefore showed steps in
brightness where one image took over from another. Where images overlap, the same location is
observed at different angles, which allows a cross-calibration. For each pair of images $(i, j)$
with at least 2000 common grid cells where $\mu, \mu_{0} > 0.3$ in both, up to \num{20000} cells
were drawn at random, and the model
\begin{equation}
\ln\frac{A_{i}}{A_{j}} = g_{i} - g_{j} +
\left[\vect{b}(\mu_{i}, \mu_{0,i}) - \vect{b}(\mu_{j}, \mu_{0,j})\right]\cdot\boldsymbol{\beta},
\qquad
\vect{b}(\mu, \mu_{0}) = \left(\bar\mu,\ \bar\mu^{2},\ \bar\mu_{0},\ \bar\mu_{0}^{2},\
\bar\mu\,\bar\mu_{0}\right),
\label{eq:selfcal}
\end{equation}
with $\bar\mu = 1 - \mu$ and $\bar\mu_{0} = 1 - \mu_{0}$, was fitted by least squares, subject to
$\sum_{i} g_{i} = 0$. Here $g_{i}$ is a gain for each image and $\boldsymbol{\beta}$ the
coefficients of an empirical correction to Equation~\eqref{eq:minnaert} shared by all images of a
set. Samples with residuals larger than four times the robust standard deviation
($1.4826$ times the median absolute deviation) were excluded and the fit repeated, four times in
all, so that clouds that moved differently from Equation~\eqref{eq:wind} do not bias the solution.
Each image is then divided by $\exp\left(g_{i} + \vect{b}\cdot\boldsymbol{\beta}\right)$.

\subsection{Temporal registration and mosaicking}

Each image is projected onto a grid of \ang{0.1} (luminance) or \ang{0.25} (colour) in
planetographic latitude and east longitude, after the clouds at each latitude have been moved to
the epoch of 20 August 1989, 08:00~UTC, by
\begin{equation}
\Delta\lambda = \frac{U(\varphi)\,\Delta t}{N(\varphi)\cos\varphi}, \qquad
N(\varphi) = \frac{a}{\sqrt{1 - e^{2}\sin^{2}\varphi}}, \qquad e^{2} = 1 - \frac{c^{2}}{a^{2}},
\label{eq:drift}
\end{equation}
where $\Delta t$ is the time from the image to the epoch and $N\cos\varphi$ the distance from the
spin axis. Cells with $\mu$ or $\mu_{0}$ below 0.25, or within 6~pixels of the edge of the valid
image area, are discarded. Each remaining cell of image $i$ receives the weight
\begin{equation}
w_{i} = \left[\frac{(\mu\,\mu_{0})^{2}\, e_{i}}{r_{i}^{2}}\right]^{3}, \qquad
e_{i} = \min\!\left(1, \frac{D_{i}}{40}\right),
\label{eq:weight}
\end{equation}
where $r_{i}$ is the resolution of the image in km per pixel and $D_{i}$ the distance of the cell
from the edge of the valid image area, in pixels. The mosaic is the weighted mean of the calibrated
albedos,
\begin{equation}
\bar{A} = \frac{\sum_{i} w_{i}\,A_{i}\,\exp\left(-g_{i} - \vect{b}_{i}\cdot\boldsymbol{\beta}\right)}
{\sum_{i} w_{i}}.
\label{eq:mosaic}
\end{equation}
The cube in Equation~\eqref{eq:weight} lets the best observation of each location dominate
without hard seams, and limits the averaging of observations made at different times, which
would otherwise double features whose motion departs from Equation~\eqref{eq:wind}.

\subsection{Colour}\label{sec:colour}

The orange, green and blue filters of the narrow-angle camera have effective wavelengths of about
591, 566 and 479~nm. They do not sample the red part of the spectrum, where absorption by methane
is strongest and which strongly affects the colour of Neptune, so they cannot provide its true
colour on their own. The mean colour was therefore taken from the reference spectrum $A(\lambda)$,
and the filter mosaics were used only to modulate it. The linear sRGB colour of a reflectance
spectrum is
\begin{equation}
\vect{c} = \sum_{\lambda} A(\lambda)\,\vect{m}(\lambda), \qquad
\vect{m}(\lambda) = \mathsf{T}\,\frac{S(\lambda)\,\bar{\vect{x}}(\lambda)}
{\sum_{\lambda'} S(\lambda')\,\bar{y}(\lambda')},
\label{eq:colour}
\end{equation}
where $S$ is the spectrum of a black body at 5778~K standing in for sunlight,
$\bar{\vect{x}} = (\bar{x}, \bar{y}, \bar{z})$ are the CIE~1931 colour-matching functions in the
analytic form of \citet{wyman2013}, and $\mathsf{T}$ is the product of the matrix from CIE~XYZ to
linear sRGB \citep{iec1999} and a Bradford chromatic adaptation \citep{lam1985} from the white
point of $S$ to D65, so that a perfect white reflector gives $(1, 1, 1)$. At each location
$\vect{x}$ of the map, the colour is
\begin{equation}
\vect{c}(\vect{x}) = \sum_{f} q_{f}(\vect{x}) \sum_{\lambda} A(\lambda)\,h_{f}(\lambda)\,
\vect{m}(\lambda),
\label{eq:modulated}
\end{equation}
where $f$ runs over the blue, green and orange filters, $q_{f}$ is the mosaic of filter $f$
divided by its mean between \ang{60}S and \ang{20}N and smoothed with a Gaussian of standard
deviation \ang{2}, and the $h_{f}$ are piecewise-linear functions that sum to one: $h_{\mathrm{B}}$
falls from 1 at 479~nm to 0 at 566~nm, $h_{\mathrm{O}}$ rises from 0 at 566~nm to 1 at 591~nm, and
$h_{\mathrm{G}} = 1 - h_{\mathrm{B}} - h_{\mathrm{O}}$. The colour images were taken three days
before the epoch, and the motions of discrete features depart from Equation~\eqref{eq:wind} over
that interval, which is why the colour is smoothed and the brightness is taken from the luminance
mosaic alone. The texture is
\begin{equation}
\vect{t}(\vect{x}) = \frac{\vect{c}(\vect{x})}{Y\!\left(\vect{c}(\vect{x})\right)}\,
Y(\bar{\vect{c}})\,\kappa\,\max\!\left(0.05,\ 1 + E\left[L'(\vect{x}) - 1\right]\right),
\label{eq:texture}
\end{equation}
where $Y(\vect{c}) = 0.2126\,c_{\mathrm{R}} + 0.7152\,c_{\mathrm{G}} + 0.0722\,c_{\mathrm{B}}$ is the
luminance, $\bar{\vect{c}}$ the colour of the unmodulated spectrum, $L'$ the luminance mosaic
divided by its mean between \ang{60}S and \ang{20}N and completed at high northern latitudes
(Section~\ref{sec:north}), and $E$ a contrast factor, equal to one for the true-colour map. The factor $\kappa$ converts the geometric albedo of the spectrum to the Minnaert
albedo of the mosaics: it is the ratio of the mean of the green mosaic between \ang{60}S and
\ang{20}N to the spectrum's albedo in a Gaussian passband of standard deviation 17~nm centred on
566~nm. The result is resampled to $4096\times2048$ pixels with a Lanczos filter in linear light
and encoded in 8-bit sRGB.

\subsection{High northern latitudes}\label{sec:north}

North of \ang{32}N all images observe the clouds close to the limb, and the data end between
\ang{37}N and \ang{47}N. Between \ang{32}N and \ang{40}N the relative luminance map $L$ is blended
into a mirror image of the southern hemisphere,
\begin{gather}
L'(\varphi, \lambda) = \left[1 - \sigma(\varphi)\right] L(\varphi, \lambda) + \sigma(\varphi)\,
\langle L(-\varphi, \cdot)\rangle_{\ang{45}},
\label{eq:fill}\\
\sigma = t^{2}(3 - 2t), \qquad t = \min\!\left(1, \max\!\left(0, \frac{\varphi - \ang{32}}{\ang{8}}
\right)\right),
\nonumber
\end{gather}
where $\langle\cdot\rangle_{\ang{45}}$ denotes a running mean over \ang{45} of longitude, applied
after a Gaussian smoothing of standard deviation \ang{0.3}, so that no individual feature appears
twice. Where $L$ is undefined, $L' = \langle L(-\varphi, \cdot)\rangle_{\ang{45}}$. The colour is
blended likewise with its mirror image. This fill is a plausible extrapolation, not data.

\section{Results}\label{sec:results}

\subsection{Pointing}

The limb fits used 67 to 221 points per image, with root-mean-square residuals between 0.37 and
0.81~pixels and a median of 0.47~pixels (Figure~\ref{fig:frame}). The scale factor is
$s = 0.9976 \pm 0.0009$ (mean and standard deviation over the 34 images), with a range of 0.9949 to
0.9996. This small, consistent departure from unity absorbs any error in the nominal pixel scale as
well as any systematic offset between the half-intensity limb and the model limb.

The summed roll-angle curves have well-defined minima (Figure~\ref{fig:twist}). The roll angle is
\ang{-6.85} on 17 August (24 images) and \ang{-6.31} on 20 August (10 images), with bootstrap
standard deviations of \ang{0.47} and \ang{0.25}. For 23 of the 34 images, the minimum of the
individual curve lies within \ang{1} of the value of its day. The result is also sensitive to
details of the processing: computing the Minnaert exponents with the biased initial disk centres
instead of the limb-fitted ones changed the roll angles by \ang{0.94} and \ang{0.19}. We therefore
regard the roll angles as uncertain by about \ang{1}.

\begin{figure}[htbp]
\centering
\includegraphics[width=0.55\textwidth]{figures/frame.jpg}
\caption{Clear-filter image C1124452, taken on 20 August 1989 at 66~km per pixel, with lines of
latitude and longitude every \ang{15} from the fitted pointing; the equator is shown in green and
the meridian \ang{0} in yellow. The disk centre and scale were fitted to the limb, the roll angle
to the cloud bands. The Great Dark Spot is to the right of the centre. The brightness is
displayed with a power of 1.8.}
\label{fig:frame}
\end{figure}

\begin{figure}[htbp]
\centering
\input{figures/twist.tex}
\caption{The roll-angle cost (Equation~\ref{eq:cost}) of each image divided by its mean (thin
lines), the mean over the images of each day (thick lines), and the roll angle adopted for each
day (dashed lines).}
\label{fig:twist}
\end{figure}

\begin{figure}[htbp]
\centering
\includegraphics[width=0.45\textwidth]{figures/blemishes.jpg}
\caption{The band-pass-filtered brightness averaged in camera coordinates over all 34 images.
The five fixed rings that are corrected are circled; the arcs are the limbs of the individual
images.}
\label{fig:blemishes}
\end{figure}

\subsection{Photometry}

The fitted Minnaert parameters (Table~\ref{tab:minnaert}) are consistent between images of the
same filter, with standard deviations of 0.003 to 0.005 in $k$, and the exponent increases towards
shorter wavelengths. The cross-calibration of the luminance set used \num{900000} pairs of samples,
of which \SI{0.4}{\percent} were rejected, and the robust standard deviation of the residuals is
\SI{2.2}{\percent}. The fitted gains lie between 0.995 and 1.007, while the shared correction ranges
from 0.942 to 1.035 between the 1st and 99th percentiles of the cells used; it is this angular
correction that removes the steps in brightness between images (Figure~\ref{fig:calibration}).

\begin{table}[htbp]
\centering
\caption{Minnaert parameters by filter: the mean exponent $k$ with its standard deviation over the
images, and the mean albedo $B$.}
\label{tab:minnaert}
\begin{tabular}{lrcc}
\toprule
Filter & Images & $k$ & $B$ \\
\midrule
Orange (591~nm) & 8 & $0.680 \pm 0.003$ & 0.397 \\
Green (566~nm) & 11 & $0.712 \pm 0.003$ & 0.505 \\
Clear & 7 & $0.752 \pm 0.003$ & 0.592 \\
Blue (479~nm) & 8 & $0.768 \pm 0.005$ & 0.628 \\
\bottomrule
\end{tabular}
\end{table}

\begin{figure}[htbp]
\centering
\includegraphics[width=0.8\textwidth]{figures/calibration.jpg}
\caption{The luminance mosaic from \ang{80}S to \ang{50}N, with its contrast stretched to show
errors, before (top) and after (bottom) the photometric cross-calibration. North is up and the
meridian \ang{0} is in the middle. Dark red marks the absence of data.}
\label{fig:calibration}
\end{figure}

\subsection{Colour}

The unmodulated spectrum gives the linear sRGB colour $\bar{\vect{c}} = (0.252, 0.471, 0.589)$, or
(138, 182, 202) in 8-bit sRGB: a pale, slightly greyish cyan blue, broadly consistent with the
colour found by \citet{irwin2024}. As a check of consistency, the ratios between the filters were
compared with those of the spectrum in Gaussian passbands of standard deviation 17~nm. The ratio
of blue to green is 1.267 in the spectrum and 1.242 in the mean Minnaert albedos of the colour
images, which agree to within \SI{2}{\percent}; the ratio of orange to green is 0.862 and 0.786. The
larger difference in orange may arise from the long-wavelength tail of the orange filter, which
extends into the methane band at 619~nm, from the simplified passbands, and from changes in the
planet between 1989 and 1993. Since the map takes only the variations of each filter about its
own mean from the Voyager images (Equation~\ref{eq:modulated}), this difference does not affect
its colour. The factor $\kappa$ in Equation~\eqref{eq:texture} is 1.155.

\subsection{The map}

The map contains data over \SI{84.8}{\percent} of the surface. It is complete in longitude up to
\ang{36.8}N and has partial coverage up to \ang{46.8}N. In true colour, the brightness varies little:
between \ang{70}S and \ang{30}N the standard deviation of $L$ is \SI{3.8}{\percent}, and
\SI{96}{\percent} of the area lies between \SI{8.2}{\percent} below and \SI{5.9}{\percent} above the
mean. The bands and the large discrete features observed by Voyager~2, including the Great Dark
Spot and its bright companion clouds, the bright feature known as the Scooter, and the second dark
spot with its bright core, are visible in the map (Figures~\ref{fig:maps}
and~\ref{fig:views}).

\begin{figure}[htbp]
\centering
\includegraphics[width=0.8\textwidth]{figures/maps.jpg}
\caption{The map in true colour ($E = 1$, top) and as distributed for the game ($E = 3$, bottom).
Equirectangular projection, north up, meridian \ang{0} in the middle, east to the right.}
\label{fig:maps}
\end{figure}

\begin{figure}[htbp]
\centering
\includegraphics[width=0.85\textwidth]{figures/views.jpg}
\caption{Orthographic views of the map with Minnaert shading ($k = 0.75$), seen from \ang{25}S and
illuminated from \ang{15} east of the viewer, rendered by the script \texttt{render.py} of the
sources and not by the game, under white light. Top: true colour. Bottom: the texture as
distributed, with its contrast stretched, without the game's tone curve. Left: the Great
Dark Spot and its companion clouds. Middle: the second dark spot and, above it, the Scooter.
Right: the hemisphere away from both.}
\label{fig:views}
\end{figure}

\section{Adaptation for StarLancer}\label{sec:game}

Two adjustments were made for the game alone; they are not part of the measured map. First, a
brightness variation of a few per cent appears featureless on a game screen, so the texture uses
$E = 3$ in Equation~\eqref{eq:texture}, in the same spirit as the contrast stretches applied to
published Voyager images of Neptune. The mean colour is unchanged by this. The original game
illuminates a planet with a single light of colour $(1, 1, 0.8)$, which would give it a yellow
cast; OpenReliant lets a mod give the planets a light of their own, and the mod makes it white, so
the texture keeps its true colour. The texture is equirectangular, with north at the top, the
meridian \ang{0} in the middle column, and east to the right, as the spheres of the mod are mapped.

Second, the limb is rendered with the part of Minnaert's law that depends on the viewing geometry.
A surface shader multiplies the texture, in linear light, by $\mu^{k-1}$ in each colour channel,
with the exponents of Table~\ref{tab:minnaert}: that of the orange filter for the red channel, for
lack of a red filter, and those of the green and blue filters. Since the law was fitted only for
$\mu \geq 0.25$, the factor is held constant closer to the limb, where it reaches about 1.5 in red
and 1.4 in blue; the centre of the disk is unchanged. As $k$ is smallest in red, the sunlit limb
becomes both brighter and paler. The game's own illumination of the planet, the cosine of the
incidence angle extended past the terminator, $\max\left(0, (\mu_{0} + 0.25)/1.25\right)$, stands
in for $\mu_{0}^{k}$. It agrees with $\mu_{0}^{0.75}$ to within \SI{1}{\percent} for
$\mu_{0} \geq 0.5$, but exceeds it by \SI{13}{\percent} at $\mu_{0} = 0.25$ and illuminates a band
of about \ang{14} beyond the terminator. Applying the full law would require the direction of
illumination within the shader, which the engine does not yet provide.

\section{Discussion}\label{sec:discussion}

At 62 to 70~km per pixel, the map is coarser than the mosaics by Jónsson
\citep{lakdawalla2014}, which were made from images taken at 3 to 4~million~km, but it is available under an open licence, and every step
of its construction can be repeated from the archived data. Its main limitations are the
following.

\paragraph{High northern latitudes.} North of \ang{40}N the map contains mirrored southern bands
(Equation~\ref{eq:fill}) rather than data, and between \ang{32}N and \ang{40}N a blend of the two.

\paragraph{Different epochs.} The luminance images span 15~hours and the colour images were taken
three days earlier. Equation~\eqref{eq:wind} describes the zonal mean motion; discrete features
drift at their own rates and change shape within hours, so their positions, and the colour around
them, are approximate. This applies to any mosaic of Neptune assembled from images taken over
hours or days.

\paragraph{Images at higher resolution.} Images from 21 and 22 August, at 30 to 53~km per pixel,
could double the resolution. We attempted to include them, but most do not contain enough of the
limb to fix their pointing, the roll angle appears to change by several degrees after 20 August,
and the clouds had changed shape within a day or two; the resulting mosaic showed curved seams and
mismatched features, and these images were left out. Including them would require the pointing of
the scan platform, or mosaics of images taken close together in time, which would then be combined.

\paragraph{Roll angle.} A single roll angle is assumed for each day, uncertain by about \ang{1}
(Section~\ref{sec:results}). An independent check against reconstructed pointing, should it
become available, would be valuable.

\paragraph{Photometry.} Minnaert's law with a shared empirical correction is a simple description
of the scattering by Neptune's hazy atmosphere. The residual scatter of \SI{2.2}{\percent} in the
cross-calibration is comparable to the contrast of many features.

\paragraph{Colour.} The reference spectrum was measured four years after the encounter, and the
filter passbands are approximated. The colour of the texture as it appears in the game is further
affected by the game's tone curve.

\paragraph{Validation.} The positions of features in the map have not been compared quantitatively
with the positions tracked by \citet{sromovsky1993}; such a comparison would provide an
independent test of the pointing and of the advection.

\section{Conclusions}\label{sec:conclusions}

We have constructed a global map of Neptune's albedo and colour from 34 Voyager~2 narrow-angle
camera images, with the pointing of each image determined from its limb and from the zonal
symmetry of the cloud bands, a photometric cross-calibration through the overlaps of the images,
and a colour derived from a ground-based spectrum. The map covers \SI{85}{\percent} of the surface
with data at 62 to 70~km per pixel. Its value lies less in the novelty of any step than in its
documentation, its reproducibility from public data, and its open licence.

\section*{Data and code availability}

The map, the scripts that produce it from the archived data, and the source of this note are
available in the \href{https://github.com/OpenReliant/openreliant-mods}{OpenReliant mods
collection}, under \texttt{sources/realistic-planets}: the list of images with their geometry and
download addresses, and Python scripts using NumPy, SciPy and ImageMagick for the download, the
pointing, the fixed-pattern correction, the mosaics and the composition, which \texttt{build.sh}
runs in order. The map and this note are released under the Creative Commons Attribution 4.0
licence, crediting OpenReliant and NASA/JPL; the scripts are released under the Mozilla Public
License 2.0.

\section*{Acknowledgements}

This work used data from the NASA Planetary Data System: the Voyager~2 images provided by the
Ring-Moon Systems Node and the spectra of E.~Karkoschka provided by the Atmospheres Node. The
Voyager images are courtesy of NASA/JPL-Caltech. The processing pipeline and this note were
developed with the assistance of an AI model (Claude, by Anthropic), under the author's
direction.

{\small
\begin{thebibliography}{99}
\bibitem[{Archinal et~al.(2018)}]{archinal2018}
Archinal, B.~A., Acton, C.~H., A'Hearn, M.~F., et~al. 2018, Report of the IAU Working Group on
Cartographic Coordinates and Rotational Elements: 2015, \textit{Celestial Mechanics and Dynamical
Astronomy}, 130, 22

\bibitem[{Danielson et~al.(1981)}]{danielson1981}
Danielson, G.~E., Kupferman, P.~N., Johnson, T.~V., \& Soderblom, L.~A. 1981, Radiometric
performance of the Voyager cameras, \textit{Journal of Geophysical Research}, 86, 8683

\bibitem[{IEC(1999)}]{iec1999}
IEC 1999, IEC 61966-2-1: Multimedia systems and equipment, Colour measurement and management,
Part 2-1: Default RGB colour space, sRGB (Geneva: International Electrotechnical Commission)

\bibitem[{Irwin et~al.(2024)}]{irwin2024}
Irwin, P.~G.~J., et~al. 2024, Modelling the seasonal cycle of Uranus's colour and magnitude, and
comparison with Neptune, \textit{Monthly Notices of the Royal Astronomical Society}, 527, 11521

\bibitem[{Karkoschka(1994)}]{karkoschka1994}
Karkoschka, E. 1994, Spectrophotometry of the Jovian planets and Titan at 300- to 1000-nm
wavelength: the methane spectrum, \textit{Icarus}, 111, 174

\bibitem[{Karkoschka(1998)}]{karkoschka1998}
Karkoschka, E. 1998, Methane spectrophotometry of the Jovian planets and Titan at 300- to 1000-nm,
ESO-J/S/N/U-SPECTROPHOTOMETER-4-V2.0, NASA Planetary Data System, doi:10.17189/2bp8-k793

\bibitem[{Lakdawalla(2014)}]{lakdawalla2014}
Lakdawalla, E. 2014, Best-ever Neptune mosaics for the 25th anniversary of Voyager 2's flyby
(mosaics by B.~Jónsson), The Planetary Society, 25 August 2014

\bibitem[{Lam(1985)}]{lam1985}
Lam, K.~M. 1985, Metamerism and colour constancy, PhD thesis, University of Bradford

\bibitem[{Limaye \& Sromovsky(1991)}]{limaye1991}
Limaye, S.~S., \& Sromovsky, L.~A. 1991, Winds of Neptune: Voyager observations of cloud motions,
\textit{Journal of Geophysical Research}, 96, 18941

\bibitem[{Minnaert(1941)}]{minnaert1941}
Minnaert, M. 1941, The reciprocity principle in lunar photometry, \textit{The Astrophysical
Journal}, 93, 403

\bibitem[{OPAL Team(2015)}]{opal}
OPAL Team 2015, Outer Planet Atmospheres Legacy (OPAL), Hubble high-level science products,
Mikulski Archive for Space Telescopes, doi:10.17909/T9G593

\bibitem[{Simon et~al.(2015)}]{simon2015}
Simon, A.~A., Wong, M.~H., \& Orton, G.~S. 2015, First results from the Hubble OPAL program:
Jupiter in 2015, \textit{The Astrophysical Journal}, 812, 55

\bibitem[{Smith et~al.(1989)}]{smith1989}
Smith, B.~A., Soderblom, L.~A., Banfield, D., et~al. 1989, Voyager 2 at Neptune: imaging science
results, \textit{Science}, 246, 1422

\bibitem[{Sromovsky et~al.(1993)}]{sromovsky1993}
Sromovsky, L.~A., Limaye, S.~S., \& Fry, P.~M. 1993, Dynamics of Neptune's major cloud features,
\textit{Icarus}, 105, 110

\bibitem[{Wyman et~al.(2013)}]{wyman2013}
Wyman, C., Sloan, P.-P., \& Shirley, P. 2013, Simple analytic approximations to the CIE XYZ color
matching functions, \textit{Journal of Computer Graphics Techniques}, 2, 1
\end{thebibliography}
}

\end{document}
